% Teaching manuscript. Edit this TeX source; Markdown is compatibility navigation only.

\StudyChapter{CS5222-Tutorial03}{Tutorial 3 · HTTP and DNS｜从报文证据走到时间计算}
\StudyTarget{CS5222-Tutorial03__tutorial-3-http-and-dns}


\StudyRef{CS5222-Chapter02}{root}{Chapter 2} · \StudyRepoRef{CS5222/course-notes/README.md}{}{课程目录} · \StudyRef{CS5222-Assignment01}{root}{Assignment 1}


依据当前 原题PDF（3页），按Problem1–4展开。教师解答共13页：Tutorial 3教师解答，第3–4张幻灯片对应Q1，Tutorial 3教师解答，第8张幻灯片对应Q2，Tutorial 3教师解答，第11–13张幻灯片对应Q3；Q4未给当前测量答案。以下保留教师结论与整理补充的区别。Q1/Q2/Q3为核心必会（中等难度，卡在字段证据与RTT条件），Q4为常规掌握的实际操作任务（需区别观察和推断）。


可选先修：\StudyRef{Foundations-NetworkBasics}{messages}{ASCII、CRLF与报文}、\StudyRef{Foundations-NetworkBasics}{timeline}{时间轴}。






\StudyHeading{1}{Problem 1. Read an HTTP request｜不知道的就明确说不知道}
\StudyTarget{CS5222-Tutorial03__q1}
\StudyTarget{CS5222-Tutorial03__problem-1-read-an-http-request}


\StudyImage{CS5222/course-notes/assets/tutorial03-01.png}


\textbf{任务 / Task:} 原题给出Wireshark看到的一段HTTP GET的ASCII内容，\StudyInline{<cr><lf>}代表实际单字节CR、LF。指出URL、HTTP版本、是否持久、浏览器主机IP、浏览器类型及该字段的用途。 / From the supplied HTTP GET bytes, identify the URL, version, connection preference, client IP, and user-agent type and purpose, citing the relevant fields.


为便于定位，下面仅摘录与答案直接相关的原字段；完整报文（包括Accept各项及原拼写）见上图：



\begin{StudyCode}
GET /cs453/index.html HTTP/1.1<cr><lf>
Host:gaia.cs.umass.edu<cr><lf>
User-Agent: Mozilla/5.0 ( Windows;U; Windows NT 5.1; en-US; rv:1.7.2)
Gecko/20040804 Netscape/7.2 (ax) <cr><lf>
Keep-Alive: 300<cr><lf>
Connection:keepalive<cr><lf><cr><lf>
\end{StudyCode}


a. \textbf{URL / 请求地址：} 按题目传统明文HTTP情境，\StudyInline{http:\allowbreak{}/\allowbreak{}/\allowbreak{}gaia.\allowbreak{}cs.\allowbreak{}umass.\allowbreak{}edu/\allowbreak{}cs453/\allowbreak{}index.\allowbreak{}html}；Host给主机，GET行给路径。若孤立看这些应用字节，并不能独立证明外层是否经过TLS，方案http来自题目情境。 / Under the exercise's plain-HTTP context, combine Host and the request target to obtain this URL.


b. \textbf{Version / 版本：} 请求行明确\StudyInline{HTTP/\allowbreak{}1.\allowbreak{}1}。 / The request line states HTTP/1.1.


c. \textbf{Connection / 连接：} 题意为持久连接。原\StudyInline{keepalive}不是规范常见的\StudyInline{keep-alive}拼写，但HTTP/1.1默认可持久，且这里未见\StudyInline{Connection:\allowbreak{} close}；不要把拼写错误当新标准token。头表达客户端意图，不单独证明服务器实际维持了多久。\href{https://www.rfc-editor.org/rfc/rfc9112}{RFC9112 §9.3} / The intended answer is persistent HTTP, consistent with HTTP/1.1 defaults; the header alone does not prove the connection's observed lifetime.


d. \textbf{Client IP / 浏览器主机IP：} 给出的HTTP头没有源IP，无法仅据此确定。\StudyInline{gaia.\allowbreak{}cs.\allowbreak{}umass.\allowbreak{}edu}是请求目标，不能当客户端。需原IP包头/捕获上下文。 / The client IP cannot be determined from these HTTP fields; inspect the enclosing IP header or capture metadata.


e. \textbf{Browser / 浏览器：} Tutorial 3教师解答，第4张幻灯片简答写Mozilla/5.0；完整User-Agent还自称Netscape/7.2，并带Mozilla/Gecko兼容信息。课堂作答可以指出教师采用的token，同时说明完整字段不只一个产品名。服务器可能据此作兼容处理、统计或内容选择；不要求每种浏览器都必须得到不同页面，也不是身份认证证据。 / The field claims Netscape 7.2 with Mozilla/Gecko compatibility tokens. It can support compatibility handling or analytics, but is self-reported and spoofable.


\textbf{独立变式 / Transfer:} \StudyInline{GET /\allowbreak{}a.\allowbreak{}png HTTP/\allowbreak{}1.\allowbreak{}1}、\StudyInline{Host:\allowbreak{} example.\allowbreak{}com}、\StudyInline{Connection:\allowbreak{} close}，能判断路径、连接意图及客户端IP吗？ / Identify the path and connection preference, and determine whether these fields reveal the client IP.



\begin{StudySolution}{答案 / Answer}


路径/a.png；请求关闭连接；不能确定客户端IP。 / Path /a.png, requested connection closure, client IP unavailable from these fields.


\end{StudySolution}






\StudyHeading{1}{Problem 2. DNS resolution｜一对箭头才是一次往返}
\StudyTarget{CS5222-Tutorial03__q2}
\StudyTarget{CS5222-Tutorial03__problem-2-dns-resolution}


\StudyImage{CS5222/course-notes/assets/tutorial03-02.png}


\textbf{条件 / Conditions:} client到local DNS的RTT为RTT\_L，local到其他DNS每次RTT为RTT\_r。a无任何缓存，按图依次访问root、TLD、authoritative；b所需记录已在local缓存。忽略额外处理与传输大小。 / The client-local RTT is RTT\_L; each local-remote DNS RTT is RTT\_r. Follow the figure's sequential root, TLD and authoritative queries without caches, then consider a local cache hit.


a 1/8合成一次client-local往返，2/3、4/5、6/7分别三次local-remote往返。总 \textbf{RTT\_L+3RTT\_r，选D}。不是8RTT，因为图上8条单向消息不能各算一次来回。


b local直接回答，有效缓存无需再查询上级，所以 \textbf{RTT\_L，选A}。本题缓存只在local；若题设改为客户端本身已有可用缓存，网络DNS往返又不同。


\textbf{English answer:} (a)RTT\_L+3RTT\_r; (b)RTT\_L. Count round trips, not individual arrows, and locate the cache precisely.


\textbf{变式 / Transfer:} RTT\_L=4ms、RTT\_r=20ms，无缓存与local命中分别多久？ / Compute both cases for these RTTs.



\begin{StudySolution}{答案 / Answer}


64ms与4ms。 / 64 ms and 4 ms. 缓存减少的是后续查询，不让client到local的往返也凭空消失。


\end{StudySolution}






\StudyHeading{1}{Problem 3. A page with eight objects｜先取目录，再知道要取哪八件东西}
\StudyTarget{CS5222-Tutorial03__q3}
\StudyTarget{CS5222-Tutorial03__problem-3-a-page-with-eight-objects}


\StudyImage{CS5222/course-notes/assets/tutorial03-03.png}


\textbf{条件 / Conditions:} 域名未缓存，顺序DNS查找代价 \(D=\sum_{i=1}^n RTT_i\)。base HTML引用同一服务器的8个很小对象；client-server RTT为RTT\textsubscript{0}，忽略发送时间。先得到HTML才能知道对象地址。 / Sequential DNS costs D=\ensuremath{\Sigma}RTT\textsubscript{i}. One base HTML file references eight small objects on the same server; server RTT is RTT\textsubscript{0} and transmission times are negligible. Object requests start after the base HTML is received.


a \textbf{Non-persistent, serial / 非持久且无并行：} 共9个对象，每个1RTT握手+1RTT请求返回，所以 \textbf{D+18RTT\textsubscript{0}，选D}。


b \textbf{Non-persistent, five parallel connections / 非持久最多5并行：} HTML先用2RTT；8对象分两批（5+3），每批各2RTT，合计 \textbf{D+6RTT\textsubscript{0}，选C}。不能把HTML也塞入尚未解析出来的8对象同一批。


c \textbf{Persistent / 持久：} 原题未明说是否pipelining；新到Tutorial 3教师解答，第13张幻灯片明确按逐个请求计算，确认采用单连接\textbf{无流水线}。结合给出的选项，一次握手+HTML一次往返+8对象各一次，共 \textbf{D+10RTT\textsubscript{0}，选B}。若另加“拿到HTML后流水线一次发完8请求”的理想条件，则D+3RTT\textsubscript{0}；这不在原选项中。因此，答案需注明采用单连接、无流水线的约定。


\textbf{English answer:} (a)D+18RTT\textsubscript{0}; (b)D+6RTT\textsubscript{0}; (c)D+10RTT\textsubscript{0} for non-pipelined persistent HTTP. Pipelining would change part c to D+3RTT\textsubscript{0} under the idealized model.


\StudyImage{CS5222/course-notes/assets/http-rounds.png}


图按握手/请求轮次计算，仅表示题目简化模型；不包括TLS、服务端处理时间、连接拥塞控制与大文件传输。


\textbf{变式 / Transfer:} 改为base HTML+6对象，非持久且最多4并行、其他条件不变，结果？ / Use six referenced objects and at most four parallel non-persistent connections.



\begin{StudySolution}{答案 / Answer}


D+2RTT\textsubscript{0}+2\ensuremath{\times}ceil(6/4)RTT\textsubscript{0}=D+6RTT\textsubscript{0}。 / D+6RTT\textsubscript{0}. 物件少两个仍需两批，所以本模型总时间没有减少。


\end{StudySolution}






\StudyHeading{1}{Problem 4. Nslookup and direct-IP browsing｜解析成功，不保证IP网址等于原站}
\StudyTarget{CS5222-Tutorial03__q4}
\StudyTarget{CS5222-Tutorial03__problem-4-nslookup-and-direct-ip-browsingip}


\textbf{原任务 / Task:} 用\StudyInline{nslookup cityu.\allowbreak{}edu.\allowbreak{}hk}找地址，再判断把IP直接输入浏览器能否得到网页。 / Resolve cityu.edu.hk and investigate whether navigating directly to its IP retrieves the site.


本机macOS于 \textbf{2026-09-23 01:19（UTC+8）} 实际查询，解析器223.5.5.5返回non-authoritative answer，两个A地址为 \textbf{45.60.199.218、45.60.197.218}。“非权威”表示答案来自递归/缓存服务器，不等于错误；多地址也不是异常。地点、时间、DNS策略变化可返回不同集合，不能把这两个IP当永久常量背下来。


直接输入IP可能打开默认页面、错误页、重定向，也可能某些服务可用；\textbf{不保证与域名访问相同}。HTTP虚拟主机依赖Host，HTTPS还涉及TLS服务器名SNI和证书名称，CDN/WAF可能按域名路由。不能从“DNS返回IP”推断“浏览器写这个IP必定显示该站”。这份2026-09-23的记录只包含DNS查询，未包含直接IP访问的结果。


\textbf{English answer:} The dated lookup returned two A records. Direct-IP navigation is not guaranteed to reproduce the hostname-based website because virtual hosting, TLS name validation and routing can depend on the hostname.



\StudyHeading{1}{来源与复习}
\StudyTarget{CS5222-Tutorial03___1}


原题p1=Problem1全部五问，p2=Problem2原图及两问，p3=Problem3三种连接模型与Problem4实际查询。理论对应Chapter2 p19–31（HTTP）、p53–67（DNS）；教师解答与本稿DNS/HTTP时间计算一致；未给出的实际网络测量与协议边界仍标为整理补充。HTTP补充标准查阅于2026-09-23。


复算脚本见\StudyRepoRef{CS5222/course-notes/tools/network_examples.py}{}{network\_examples.py}。关联\href{https://crazyshout.github.io/micro-course/cards.html\#CS5222-N231}{现有原题卡N231、N232、N233}时，先能解释字段证据和计时条件，再用Markji巩固。
